% ============================================================
%  Chapter 2 — The Quadratic Sequence
% ============================================================
\section{The Quadratic Sequence}\label{sec:sequence}

\subsection{Definition and Origin}

The sequence arises from the classical summation formula for squares.
The sum of the first $n$ squares is
$S_n = \sum_{i=1}^{n} i^2 = \frac{n(n+1)(2n+1)}{6}$,
and its arithmetic mean
\begin{equation}\label{eq:k-explicit}
  k(n) = \frac{S_n}{n} = \frac{(n+1)(2n+1)}{6}
       = \frac{2n^2 + 3n + 1}{6}
\end{equation}
is a quadratic polynomial in $n$ with leading coefficient $\frac{1}{3}$.

\begin{remark}[OEIS entry]
The integer values of $k(n)$ are listed in the OEIS under
\href{https://oeis.org/A164576}{A164576}~\cite{OEIS} as
\emph{``Integer averages of the set of the first positive squares up to
some $n^2$''}.
\end{remark}

\subsection{Integrality Condition}

\begin{theorem}[Congruence condition]\label{thm:congruence}
The sequence $k(n)$ takes integer values if and only if
\begin{equation}\label{eq:cong}
  n \equiv 1, 5 \pmod{6}.
\end{equation}
\end{theorem}

\begin{proof}
We examine $2n^2 + 3n + 1 \pmod{6}$ for all residue classes:

\begin{center}
\begin{tabular}{ccccc}
\toprule
$n \!\pmod{6}$ & $2n^2 \!\pmod{6}$ & $3n \!\pmod{6}$ & $2n^2\!+\!3n\!+\!1 \!\pmod{6}$ & $6 \mid ?$ \\
\midrule
0 & 0 & 0 & 1 & No \\
1 & 2 & 3 & 0 & Yes \\
2 & 2 & 0 & 3 & No \\
3 & 0 & 3 & 4 & No \\
4 & 2 & 0 & 3 & No \\
5 & 2 & 3 & 0 & Yes \\
\bottomrule
\end{tabular}
\end{center}
Divisibility by~$6$ occurs precisely for $n \equiv 1$ and $n \equiv 5$.
\end{proof}

The first integer values are:

\begin{center}
\begin{tabular}{c*{10}{c}}
\toprule
$n$ & 1 & 5 & 7 & 11 & 13 & 17 & 19 & 23 & 25 & 29 \\
\midrule
$k(n)$ & 1 & 11 & 20 & 46 & 63 & 105 & 130 & 188 & 221 & 295 \\
\bottomrule
\end{tabular}
\end{center}

\subsection{Difference Structure: Long and Short Steps}\label{subsec:steps}

The integer-valued indices follow a periodic pattern. Within each
period $m = 1, 2, 3, \ldots$, there is a step of width $\Delta n = 4$
from $n = 6m-5$ to $n = 6m-1$, followed by a step of width
$\Delta n = 2$ from $n = 6m-1$ to $n = 6m+1$.

\begin{definition}[Long and short steps]\label{def:steps}
For the $m$-th period ($m \geq 1$) we define:
\begin{align}
  \delta^{(4)}_m &\coloneqq k(6m-1) - k(6m-5)
  \quad\text{(\emph{long step},\; $\Delta n = 4$)},
  \label{eq:long}\\[4pt]
  \delta^{(2)}_m &\coloneqq k(6m+1) - k(6m-1)
  \quad\text{(\emph{short step},\; $\Delta n = 2$)}.
  \label{eq:short}
\end{align}
The superscript indicates the index difference $\Delta n$.
\end{definition}

\begin{theorem}[Step-size formulae]\label{thm:step-formulae}
For all $m \geq 1$:
\begin{equation}\label{eq:step-formulae}
  \delta^{(4)}_m = 16m - 6,
  \qquad
  \delta^{(2)}_m = 8m + 1.
\end{equation}
\end{theorem}

\begin{proof}
Using $k(n) = \frac{1}{6}(2n^2+3n+1)$ and the difference-of-squares identity:
\begin{align*}
  k(6m\!-\!1) - k(6m\!-\!5)
  &= \frac{1}{6}\bigl[2\cdot(12m\!-\!6)\cdot 4 + 3\cdot 4\bigr]
  = \frac{96m-36}{6} = 16m-6,\\[4pt]
  k(6m\!+\!1) - k(6m\!-\!1)
  &= \frac{1}{6}\bigl[2\cdot 12m\cdot 2 + 3\cdot 2\bigr]
  = \frac{48m+6}{6} = 8m+1.
  \qedhere
\end{align*}
\end{proof}

The first differences are thus:
\begin{center}
\begin{tabular}{c*{8}{c}}
\toprule
$\delta_k$ & 10 & 9 & 26 & 17 & 42 & 25 & 58 & 33 \\
$\Delta n$ & 4 & 2 & 4 & 2 & 4 & 2 & 4 & 2 \\
\bottomrule
\end{tabular}
\end{center}

\begin{remark}[Pairwise sums form an arithmetic progression]
Each long--short pair sums to
$U_m = \delta^{(4)}_m + \delta^{(2)}_m = (16m-6)+(8m+1) = 24m-5$,
with constant increment $\Delta U = 24$. This arithmetic structure is a
direct algebraic consequence of Theorem~\ref{thm:step-formulae} and will
reappear as the circumference of the $m$-th spiral winding in
Section~\ref{sec:helix}.
\end{remark}

\subsection{Prime Number Property}\label{subsec:primes-mod6}

\begin{theorem}[Primes and the congruence classes]\label{thm:primes-mod6}
All primes $p > 3$ satisfy $p \equiv 1$ or $5 \pmod{6}$.
\end{theorem}

\begin{proof}
Every prime $p > 3$ is coprime to $6$, since it is divisible neither by~$2$
nor by~$3$. The residue classes coprime to~$6$ are precisely $1$ and
$5$~\cite[Ch.\,1]{IrelandRosen1990}.
\end{proof}

Thus, the indices at which $k(n)$ is an integer are exactly the residue
classes that contain all primes greater than~$3$. This coincidence is the arithmetic foundation of the prime number
corollary in Section~\ref{sec:primes}.
